% Generated draft: requires mathematical and source review.
\documentclass[11pt]{article}
\usepackage[margin=25mm]{geometry}
\usepackage{fontspec}
\setmainfont{Latin Modern Roman}
\usepackage{amsmath,amssymb,mathtools,unicode-math}
\setmathfont{Latin Modern Math}
\usepackage{enumitem,needspace,xurl,hyperref}
\setlength{\parindent}{0pt}
\setlength{\parskip}{5pt}
\setlength{\emergencystretch}{3em}
\newcommand{\N}{\mathbb{N}}
\newcommand{\Z}{\mathbb{Z}}
\newcommand{\Q}{\mathbb{Q}}
\newcommand{\R}{\mathbb{R}}
\newcommand{\C}{\mathbb{C}}
\newcommand{\F}{\mathbb{F}}
\newcommand{\A}{\mathbb{A}}
\newcommand{\PP}{\mathbb{P}}
\newcommand{\E}{\mathbb{E}}
\newcommand{\eps}{\varepsilon}
\newcommand{\dd}{\,\mathrm{d}}
\begin{document}
\section*{UnsolvedMath ID: 30004677}
\textbf{Source problem code:} catalogue entry
\subsection{Definitions and mathematical statement}
The supplied record asks two questions about degrees and the skip operator, but does not specify their underlying definitions. To state the questions relative to these missing data, let \((\mathcal D,\leq)\) denote the intended partially ordered degree structure, let \(\mathcal A\subseteq\mathcal D\) be its arithmetical degrees, let \(\mathcal T\subseteq\mathcal D\) be its total degrees, and let \(\mathcal H\subseteq\mathcal D\) be the class called HYP degrees in the source. Write \(S:\mathcal D\to\mathcal D\) for the intended skip operator. These symbols are placeholders for source-specific definitions, not definitions of the underlying computability-theoretic notions.

With the iteration convention \(S^0(\mathbf a)=\mathbf a\) and \(S^{n+1}(\mathbf a)=S(S^n(\mathbf a))\), the first question, if the source includes the zeroth iterate, is whether
\[
\exists\mathbf a\in\mathcal A\;\forall n\in\{0,1,2,\ldots\},\qquad S^n(\mathbf a)\not\leq S^{n+1}(\mathbf a).
\]
The source presents this as equivalent to
\[
\exists\mathbf a\in\mathcal A\;\forall n\in\{0,1,2,\ldots\},\qquad S^n(\mathbf a)\notin\mathcal T.
\]
This equivalence is supplied source content; its justification depends on the missing definitions. The record does not explicitly specify whether its phrase ``for every \(n\)'' includes \(n=0\). If only positive iterates are intended, both displayed quantifiers must instead range over \(\{1,2,\ldots\}\).

For a reading in which a skip \(2\)-cycle has two distinct elements, put
\[
\mathcal C_2=\{(\mathbf b,\mathbf c)\in\mathcal D^2:\mathbf b\ne\mathbf c,\ S(\mathbf b)=\mathbf c,\ S(\mathbf c)=\mathbf b\}.
\]
If ``bounded by all skip \(2\)-cycles'' means lying below both elements of every such cycle, the second question is whether
\[
\forall\mathbf d\in\mathcal D,\qquad
\mathbf d\in\mathcal H
\ \Longleftrightarrow\
\forall(\mathbf b,\mathbf c)\in\mathcal C_2,\quad
(\mathbf d\leq\mathbf b\ \text{and}\ \mathbf d\leq\mathbf c).
\]
This is a conditional transcription: the supplied record does not explicitly settle whether fixed points count as \(2\)-cycles, whether both cycle elements must bound \(\mathbf d\), or the ambient domain of the characterization. The degree structure, reducibility, skip construction, and meanings of arithmetical, total, and HYP degrees require confirmation from the cited sources before these formulas constitute a fully specified statement.
\subsection{Short English statement}
The first question asks for one arithmetical degree whose successive skips always fail to lie below the next skip; the source calls this equivalent to every iterate being nontotal. The second asks whether HYP degrees are exactly the common lower bounds of all skip \(2\)-cycles. Several underlying definitions and conventions remain unspecified. Supplied dated curation reports partial progress, including an affirmative answer to the second question attributed to Goh; this draft does not independently verify that assessment.
\subsection{Sources}
\begin{itemize}
\item Local UnsolvedMath registry, permanent ID \(30004677\), title ``Zigzagging Skips of Arithmetical Degrees,'' source problem code \texttt{catalogue entry}, in the Oberwolfach Reports Open Problems collection. The supplied proposed year is 2021; no individual proposer is identified.
\item Supplied report citation: \emph{Computability Theory (2022)}, Oberwolfach Reports, DOI \texttt{10.4171/owr/2021/21}, pp.\ 1149--1190. \url{https://doi.org/10.4171/owr/2021/21}. The supplied background describes Oberwolfach Reports as documenting mathematical workshops and their open-problem discussions.
\item Andrews et al., \emph{On cototality and the skip operator}. \url{https://people.math.wisc.edu/~slempp/papers/cototal.pdf}.
\item Soskova, \emph{The e-verse}. \url{https://people.math.wisc.edu/~soskova/talks/WDCM2021Soskova.pdf}.
\item \emph{Oberwolfach Report 21/2021}. \url{https://ems.press/content/serial-article-files/46899}.
\item Literature assessment dated \texttt{2026-08-22}: the supplied curation labels the problem partially solved, reports that skip \(2\)-cycles were constructed and that Goh answered the HYP-degree characterization affirmatively, and states that those results do not settle the separate existence question for a single degree with the required zig-zag behavior. This is supplied dated curation derived from the OWR record, not an independently verified literature assessment. Its time-sensitive status and cited claims require independent verification.
\end{itemize}
\end{document}
